\documentclass[12pt]{article}

% a template that a friend gave, it's worked well enough for me
% i have added some packages and stuff that have proved useful

\usepackage{fancyhdr}
\usepackage{tipa}
\usepackage{fontspec}
\usepackage{amsfonts}
\usepackage{enumitem}
\usepackage[margin=1in]{geometry}
\usepackage{graphicx}
\usepackage{float}
\usepackage{amsmath}
\usepackage{braket}
\usepackage{amssymb}
\usepackage{booktabs}
\usepackage{hyperref}
\usepackage{mathtools}
\usepackage{xcolor}
\usepackage{float}
\usepackage{algpseudocodex}
\usepackage{titlesec}
\usepackage{bbm}

\pagestyle{fancy}
\fancyhf{} % sets both header and footer to nothing
\lhead{Kevin Sheng}
\setmainfont{Comic Neue}
\renewcommand{\headrulewidth}{1pt}
\setlength{\headheight}{0.75in}
\setlength{\oddsidemargin}{0in}
\setlength{\evensidemargin}{0in}
\setlength{\voffset}{-.5in}
\setlength{\headsep}{10pt}
\setlength{\textwidth}{6.5in}
\setlength{\headwidth}{6.5in}
\setlength{\textheight}{8in}
\renewcommand{\headrulewidth}{0.5pt}
\renewcommand{\footrulewidth}{0.3pt}
\setlength{\textwidth}{6.5in}
\usepackage{setspace}
\usepackage{multicol}
\usepackage{float}
\setlength{\columnsep}{1cm}
\setlength\parindent{24pt}
\usepackage [english]{babel}
\usepackage [autostyle, english = american]{csquotes}
\MakeOuterQuote{"}

\setlength{\parskip}{6pt}
\setlength{\parindent}{0pt}

\titlespacing\section{0pt}{12pt plus 4pt minus 2pt}{0pt plus 2pt minus 2pt}
\titlespacing\subsection{0pt}{12pt plus 4pt minus 2pt}{0pt plus 2pt minus 2pt}
\titlespacing\subsubsection{0pt}{12pt plus 4pt minus 2pt}{0pt plus 2pt minus 2pt}

\hypersetup{colorlinks=true, urlcolor=blue}

\newcommand{\correction}[1]{\textcolor{red}{#1}}

\usepackage{fancyvrb}

\rhead{Math 168}

\DeclareMathOperator{\Fr}{Fr}
\DeclareMathOperator{\Tr}{Tr}
\newcommand{\lra}{\xLeftrightarrow}
\newcommand{\ra}{\xRightarrow}

\setlist[enumerate,1]{label=(\alph*)}

\begin{document}

\section{Newman 12.3}

\begin{enumerate}
      \item $\braket{x}_{edge}$ just seems to be the average value at all edge ends.
            Since we know the degree of each node, the total value of all edge
            ends combined is $\sum_{i} k_i x_i$, and since there's $\frac{1}{2m}$ edge ends,
            $\braket{x}_{edge}=\frac{1}{2m} \sum_{i} k_i x_i$.

      \item Consider any edge emerging out of node $i$.
            As derived in class, the chance this edge hits a specific node that has degree $k$ is around $\frac{k}{2m}$.
            If we sum over all possible nodes and average their values, we have
            the average value of a node's neighbor as
            \[\sum_{i=1}^{n} x_i \cdot \frac{k_i}{2m} = \braket{x}_{edge}\]
            If we subtract $\braket{x}=\frac{1}{n} \sum_{i} x_i$, we have
            \begin{align*}
                  \braket{x}_{edge} - \braket{x}
                   & = \frac{1}{2m} \sum_{i=1}^{n} x_i k_i - \frac{1}{n} \sum_{i=1}^{n} x_i                       \\
                   & = \frac{1}{n\braket{k}} \sum_{i=1}^{n} x_i k_i - \frac{1}{n} \sum_{i=1}^{n} x_i              \\
                   & = \frac{\frac{1}{n}\sum_{i=1}^{n} x_i k_i - \sum_{i=1}^{n} x_i \cdot \braket{k}}{\braket{k}} \\
                   & = \frac{\braket{kx} - \braket{x}\braket{k}}{\braket{k}}                                      \\
                   & = \frac{\text{cov}(k, x)}{k}\quad\square
            \end{align*}

      \item Since all edges are constructed independently, just because
            I'm friends with a neighbor doesn't make them any more likely to
            have more friends than I do.
\end{enumerate}

\pagebreak

\section{ER Graphs' Local Trees}

Assume that $\braket{k}=p(n-1)$ is constant as $n \to \infty$.

I guess we should consider all nodes that are at most $d$ away from a fixed node $v$.
There's approximately $\braket{k}^i$ nodes that are distance exactly $i$ from $v$,
so in total our local connected component should have around
\[s=\sum_{i=0}^{d} \braket{k}^i = \frac{\braket{k}^{d+1}-1}{\braket{k}-1}\]
nodes in total.

We're looking for a local tree, so it should have exactly $s-1$ edges.
Since there's $\binom{s}{2}$ edge candidates, the chance of no other
edges appearing besides the ones that we took to get to all the nodes is $(1-p)^{\binom{s}{2}-s+1}$.

Though the exponent is large, it actually doesn't really matter.
As $n \to \infty$, $p \to 0$.
Since our derived expression for $s$ has no dependence on $n$, the exponent
stays the same while the base slowly goes to $1$.
Thus, holding $\braket{k}$ constant while taking $n \to \infty$ means
that ER graphs are probably gonna be locally tree-like for some $d$.

\pagebreak

\section{Newman 12.7}

\begin{enumerate}
      \item It STP that $p_k=\frac{g^{(k)}(0)}{k!}$.

            The proof is by induction on the expression for the RHS.
            I propose
            \[\frac{g^{(k)}(z)}{k!}=\binom{n}{k}p^k(pz+1-p)^{n-k}\]
            For $k=0$, we have $\frac{g(z)}{0!}=(pz+1-p)^n$, which aligns with the proposed formula.

            Assuming truth for $k$, we see have
            \begin{align*}
                  \frac{g^{(k+1)}(z)}{(k+1)!}
                   & = \frac{1}{k+1} \cdot \frac{d}{dz} \frac{g^{(k)}(z)}{k!}                       \\
                   & = \frac{1}{k+1} \cdot \frac{d}{dz} \left(\binom{n}{k}p^k(pz+1-p)^{n-k}\right)  \\
                   & = \frac{1}{k+1} \cdot \binom{n}{k} \cdot p^k \cdot \frac{d}{dz} (pz+1-p)^{n-k} \\
                   & = \frac{1}{k+1} \cdot \binom{n}{k} \cdot p^k \cdot (n-k)(pz+1-p)^{n-k-1}p      \\
                   & = \frac{\binom{n}{k} (n-k)}{k+1} \cdot p^{k+1}(pz+1-p)^{n-k-1}                 \\
                   & = \binom{n}{k+1}p^{k+1}(pz+1-p)^{n-k-1}
            \end{align*}
            Evaluating all of these formulas at $0$ indeed gives $p_k=\frac{g^{(k)}(0)}{k!}$. $\square$

      \item The first moment is
            \[g'(1)=\binom{n}{1}p^1(p \cdot 1+1-p)^{n-1}=np\]
            and the second moment is

      \item For a single $k$, assuming $0 \le k \le 2n$, we have
\end{enumerate}

\section{Sanity Simulations}

See the attached Python notebook.

\end{document}
